%%
%% part2-book-components/ch13-margin-notes.tex
%%
%% Chapter 13: Margin Notes — RTL margin notes and
%% margin figures in the wide outer margin.
%%

\chapter{یادداشت حاشیه}
\label{chap:margin-notes}

\section{چرا این موضوع مهم است؟}
\label{sec:margin-notes-why}

یکی از ویژگی‌های متمایزکننده‌ی کتاب‌های درسی
حرفه‌ای، استفاده از \textbf{حاشیه‌ی بیرونی پهن}
برای یادداشت‌های توضیحی، شکل‌های کوچک و
فرمول‌های مکمل است. این طراحی، از کتاب‌های
\lr{Boyer} و \lr{Stewart} الهام گرفته شده
است.

\lr{MatinBook} یک حاشیه‌ی بیرونی پهن (\pnum{۵.۹}
سانتی‌متر) ارائه می‌دهد که مخصوص این کار طراحی
شده است. در این فصل، با دو دستور \cmd{mnote}
و \cmd{marginfig} آشنا می‌شوید.

\mnote{حاشیه‌ی پهن در \lr{frontmatter} و \lr{backmatter} حذف می‌شود و فقط در \lr{mainmatter} فعال است.}

\section{دستور \cmd{mnote}}
\label{sec:margin-notes-mnote}

دستور \cmd{mnote} برای افزودن یک یادداشت متنی
ساده در حاشیه استفاده می‌شود:

\begin{latin}
\begin{minted}{latex}
|\pc{این یک متن اصلی است.}|\mnote{|\pc{این یک یادداشت حاشیه‌ای است.}|}%
|\pc{و این ادامه‌ی متن اصلی است.}|%
\end{minted}
\end{latin}

خروجی:

این یک متن اصلی است.\mnote{این یک یادداشت حاشیه‌ای است.}
و این ادامه‌ی متن اصلی است.

\mnote{یادداشت‌های حاشیه‌ای باید کوتاه و مکمل متن اصلی باشند. طول ایده‌آل: ۱ تا ۳ خط.}

\subsection{ویژگی‌های \cmd{mnote}}
\label{sec:margin-notes-mnote-features}

دستور \cmd{mnote} ویژگی‌های زیر را دارد:

\begin{itemize}
    \item \textbf{بدون شماره:} یادداشت‌ها
    شماره‌گذاری نمی‌شوند (برخلاف پانویس).
    \item \textbf{راست‌به‌چپ:} متن یادداشت
    به‌طور خودکار راست‌به‌چپ است.
    \item \textbf{حاشیه‌ی بیرونی:} در صفحات فرد،
    در سمت راست؛ در صفحات زوج، در سمت چپ.
    \item \textbf{اندازه‌ی کوچک:} متن یادداشت با
    \cmd{footnotesize} نمایش داده می‌شود.
\end{itemize}

\mnote{یادداشت‌های حاشیه‌ای با \lr{\textbackslash marginpar} پیاده‌سازی شده‌اند که بخشی از هسته‌ی \lr{LaTeX} است.}

\section{دستور \cmd{marginfig}}
\label{sec:margin-notes-marginfig}

دستور \cmd{marginfig} برای افزودن یک شکل کوچک
با کپشن در حاشیه استفاده می‌شود:

\begin{latin}
\begin{minted}{latex}
|\pc{این یک متن اصلی است.}|%
\marginfig{figures/plot.png}{|\pc{نمودار نمونه}|}%
\end{minted}
\end{latin}

این دستور دو آرگومان می‌گیرد:

\begin{enumerate}
    \item \textbf{نام فایل:} مسیر فایل تصویر.
    \item \textbf{کپشن:} متن توضیح زیر شکل.
\end{enumerate}

\mnote{عرض شکل در حاشیه، به‌طور خودکار با \lr{\textbackslash marginparwidth} تنظیم می‌شود که \lr{۳.۵} سانتی‌متر است.}

\section{محدودیت‌ها}
\label{sec:margin-notes-limitations}

یادداشت‌های حاشیه‌ای چند محدودیت دارند:

\subsection{عدم استفاده در \env{tcolorbox}}
\label{sec:margin-notes-limitation-tcolorbox}

دستور \cmd{mnote} \textbf{نمی‌تواند} داخل
\env{tcolorbox} (مثل \env{theorem}، \env{definition}
و...) استفاده شود، چون \cmd{marginpar} در
محیط‌های شناور کار نمی‌کند.

\textbf{راه‌حل:} یادداشت را بیرون از جعبه قرار دهید.

\subsection{عدم استفاده در \env{figure}}
\label{sec:margin-notes-limitation-figure}

دستور \cmd{mnote} \textbf{نمی‌تواند} داخل
\env{figure} استفاده شود.

\textbf{راه‌حل:} یادداشت را در متن اصلی، بیرون
از \env{figure}، قرار دهید.

\subsection{نیاز به هندسه‌ی \lr{mainmatter}}
\label{sec:margin-notes-limitation-geometry}

یادداشت‌های حاشیه‌ای فقط در \lr{mainmatter} کار
می‌کنند، چون در \lr{frontmatter} و \lr{backmatter}
حاشیه‌ی پهن حذف می‌شود.

\textbf{راه‌حل:} از \cmd{mainmattergeometry} برای
فعال‌سازی حاشیه‌ی پهن استفاده کنید.

\subsection{کامپایل چندباره}
\label{sec:margin-notes-limitation-compile}

یادداشت‌های حاشیه‌ای برای موقعیت‌یابی پایدار،
حداقل \textbf{۲ بار کامپایل} نیاز دارند.

\textbf{راه‌حل:} چرخه‌ی کامل کامپایل را اجرا
کنید.

\mnote{اگر یادداشت‌ها در جای اشتباه ظاهر شدند، یک بار دیگر کامپایل کنید.}

\section{جهت متن در صفحات زوج و فرد}
\label{sec:margin-notes-direction}

یادداشت‌های حاشیه‌ای در \lr{MatinBook} در
هر دو نوع صفحه (فرد و زوج) راست‌به‌چپ هستند.
این کار با استفاده از \cmd{RL} از \lr{xepersian}
انجام می‌شود.

جدول \cref{tab:margin-direction} رفتار یادداشت‌ها
را در صفحات مختلف نشان می‌دهد.

\begin{matintable}{رفتار یادداشت‌های حاشیه‌ای}{tab:margin-direction}
\begin{tblr}{
    width = \textwidth,
    colspec = {l X[c] X[c]},
    hlines, vlines,
    colsep = 8pt,
    rowsep = 4pt,
    row{1} = {font=\bfseries},
}
نوع صفحه & سمت حاشیه & جهت متن \\
فرد (راست) & راست & راست‌به‌چپ \\
زوج (چپ) & چپ & راست‌به‌چپ \\
\end{tblr}
\end{matintable}

\mnote{در \lr{MatinBook}، از \lr{\textbackslash reversemarginpar} استفاده نمی‌شود، چون جهت را دو بار معکوس می‌کند.}

\section{مثال‌های کاربردی}
\label{sec:margin-notes-examples}

\subsection{یادداشت توضیحی}
\label{sec:margin-notes-example-explanation}

یادداشت‌ها برای توضیحات مکمل مناسب هستند:

\begin{latin}
\begin{minted}{latex}
|\pc{الگوریتم دایکسترا برای یافتن کوتاه‌ترین مسیر در گراف‌های وزن‌دار استفاده می‌شود.}|%
\mnote{|\pc{پیچیدگی این الگوریتم}|\lr{$O((V+E) \log V)$}|\pc{است.}|}%
\end{minted}
\end{latin}

خروجی:

الگوریتم دایکسترا برای یافتن کوتاه‌ترین مسیر در
گراف‌های وزن‌دار استفاده می‌شود.\mnote{پیچیدگی این الگوریتم \lr{$O((V+E) \log V)$} است.}

\subsection{یادداشت هشدار}
\label{sec:margin-notes-example-warning}

یادداشت‌ها برای هشدارها نیز مناسب هستند:

\begin{latin}
\begin{minted}{latex}
|\pc{تقسیم بر صفر تعریف نشده است.}|%
\mnote{|\pc{هنگام ساده‌سازی عبارات کسری، همواره فرض کنید مخرج غیرصفر است.}|}%
\end{minted}
\end{latin}

خروجی:

تقسیم بر صفر تعریف نشده است.\mnote{هنگام ساده‌سازی عبارات کسری، همواره فرض کنید مخرج غیرصفر است.}

\subsection{یادداشت ارجاع}
\label{sec:margin-notes-example-reference}

یادداشت‌ها برای ارجاع به منابع نیز مناسب هستند:

\begin{latin}
\begin{minted}{latex}
|\pc{این قضیه اولین بار توسط اویلر اثبات شد.}|%
\mnote{|\pc{برای مطالعه‌ی بیشتر، به}|\cite{knuth1984}|\pc{مراجعه کنید.}|}%
\end{minted}
\end{latin}

\section{شخصی‌سازی}
\label{sec:margin-notes-customization}

برای شخصی‌سازی یادداشت‌های حاشیه‌ای، می‌توانید
پارامترهای زیر را در فایل \file{mb-layout.sty}
تغییر دهید:

\begin{matintable}{پارامترهای قابل تنظیم یادداشت حاشیه}{tab:margin-params}
\begin{tblr}{
    width = \textwidth,
    colspec = {l X[l] X[c]},
    hlines, vlines,
    colsep = 6pt,
    rowsep = 4pt,
    row{1} = {font=\bfseries},
}
پارامتر & توضیح & مقدار پیش‌فرض \\
\lr{marginparwidth} & عرض ناحیه‌ی یادداشت & \pnum{۳.۵} سانتی‌متر \\
\lr{marginparsep} & فاصله از متن اصلی & \pnum{۰.۶} سانتی‌متر \\
\lr{outer} & حاشیه‌ی بیرونی کل & \pnum{۵.۹} سانتی‌متر \\
\end{tblr}
\end{matintable}

\mnote{تغییر این پارامترها، بر چیدمان کل کتاب تأثیر می‌گذارد. با احتیاط عمل کنید.}

\section{تفاوت با پانویس}
\label{sec:margin-notes-vs-footnote}

یادداشت‌های حاشیه‌ای با پانویس تفاوت‌هایی
دارند. جدول \cref{tab:margin-vs-footnote} این
تفاوت‌ها را نشان می‌دهد.

\begin{matintable}{تفاوت یادداشت حاشیه و پانویس}{tab:margin-vs-footnote}
\begin{tblr}{
    width = \textwidth,
    colspec = {l X[c] X[c]},
    hlines, vlines,
    colsep = 6pt,
    rowsep = 4pt,
    row{1} = {font=\bfseries},
}
ویژگی & یادداشت حاشیه & پانویس \\
موقعیت & حاشیه‌ی صفحه & پایین صفحه \\
شماره‌گذاری & ندارد & دارد \\
طول & کوتاه (۱-۳ خط) & متوسط (۱-۵ خط) \\
کاربرد & توضیح مکمل & ارجاع یا توضیح \\
دستور & \cmd{mnote} & \cmd{footnote} \\
\end{tblr}
\end{matintable}

\mnote{از یادداشت حاشیه برای توضیحات کوتاه و مکمل استفاده کنید، و از پانویس برای ارجاعات و توضیحات بلندتر.}

\section{خطاهای رایج}
\label{sec:margin-notes-mistakes}

\subsection{استفاده در \env{tcolorbox}}
\label{sec:margin-notes-mistake-tcolorbox}

\textbf{مشکل:} \cmd{mnote} داخل \env{tcolorbox}
(مثل \env{theorem}، \env{definition}) کار نمی‌کند
و خطا می‌دهد.

\textbf{راه‌حل:} یادداشت را بیرون از جعبه قرار
دهید.

\subsection{استفاده در \env{figure}}
\label{sec:margin-notes-mistake-figure}

\textbf{مشکل:} \cmd{mnote} داخل \env{figure} کار
نمی‌کند.

\textbf{راه‌حل:} یادداشت را در متن اصلی قرار
دهید.

\subsection{یادداشت بسیار طولانی}
\label{sec:margin-notes-mistake-long}

\textbf{مشکل:} اگر یادداشت بیش از ۳ خط باشد،
ممکن است از حاشیه سرریز کند یا به صفحه‌ی بعد
برود.

\textbf{راه‌حل:} یادداشت‌ها را کوتاه نگه دارید
(حداکثر ۳ خط). برای توضیحات بلندتر، از متن اصلی
یا پانویس استفاده کنید.

\subsection{فراموش کردن هندسه‌ی \lr{mainmatter}}
\label{sec:margin-notes-mistake-geometry}

\textbf{مشکل:} اگر \cmd{mainmattergeometry} را
فراخوانی نکنید، حاشیه‌ی پهن وجود ندارد و
یادداشت‌ها در جای اشتباه ظاهر می‌شوند.

\textbf{راه‌حل:} همیشه پس از \cmd{mainmatter}،
دستور \cmd{mainmattergeometry} را فراخوانی کنید.

\section{تمرین‌ها}
\label{sec:margin-notes-exercises}

\begin{exercise}
    یک متن ساده با دو یادداشت حاشیه‌ای بنویسید.
    آیا یادداشت‌ها در حاشیه‌ی درست ظاهر
    می‌شوند؟
    \label{exc:margin-notes-basic}
\end{exercise}

\begin{exercise}
    یک یادداشت هشدار در حاشیه بنویسید که
    به یک بخش دیگر کتاب ارجاع دهد.
    \label{exc:margin-notes-warning}
\end{exercise}

\begin{exercise}
    یک \env{theorem} بنویسید و یک یادداشت
    حاشیه‌ای در بیرون آن (نه داخلش) قرار
    دهید.
    \label{exc:margin-notes-outside}
\end{exercise}

\begin{exercise}
    یک یادداشت بسیار طولانی بنویسید و ببینید
    آیا سرریز می‌کند. سپس آن را کوتاه کنید.
    \label{exc:margin-notes-long}
\end{exercise}

\section{خلاصه}
\label{sec:margin-notes-summary}

در این فصل، با یادداشت حاشیه‌ای در \lr{MatinBook}
آشنا شدیم:

\begin{itemize}
    \item \lr{MatinBook} یک حاشیه‌ی بیرونی پهن
    (\pnum{۵.۹} سانتی‌متر) برای یادداشت‌ها ارائه
    می‌دهد.

    \item دو دستور برای یادداشت:
    \cmd{mnote} (متنی) و \cmd{marginfig} (شکل).

    \item یادداشت‌ها بدون شماره هستند و جهت
    راست‌به‌چپ دارند.

    \item یادداشت‌ها در حاشیه‌ی بیرونی قرار
    می‌گیرند (راست در صفحات فرد، چپ در
    صفحات زوج).

    \item محدودیت‌ها: عدم استفاده در
    \env{tcolorbox} و \env{figure}، نیاز به
    \cmd{mainmattergeometry}، و کامپایل چندباره.

    \item یادداشت‌ها با پانویس تفاوت دارند:
    بدون شماره، کوتاه‌تر، و در حاشیه.

    \item خطاهای رایج شامل استفاده در
    \env{tcolorbox}، استفاده در \env{figure}،
    یادداشت بسیار طولانی، و فراموش کردن
    \cmd{mainmattergeometry} است.
\end{itemize}

در فصل بعدی (\cref{chap:bibliography})، با
منابع و کتاب‌نامه در \lr{MatinBook} آشنا
می‌شوید و یاد می‌گیرید که چگونه به منابع
ارجاع دهید.